\section{Conclusions}
\label{sec:conclusions}

Wrapping a value-at-risk model in an adaptive conformal update of its tail level improves its coverage in market stress. On ten of the eighteen models of the ladder and three stress windows, the wrapped models pass the Kupiec and conditional coverage tests in \tenPassAci{} of \tenRuns{} model-level-window cells against \tenPassUnw{} unwrapped. The prediction stated before the runs holds: HS-250, GARCH-t and filtered historical simulation each gain at least one stress window at both levels (\hypCellsHold{} of \hypCells{} cells), the worst change of the mean score is \hypWorstFz\%, below the 5\% tolerance, and the independence test is not failed in the crisis window (p-values \indMin{} to \indMax{}). On generators of known truth the wrapped pass rate is higher in all \genPredANum{} cells with a shift.

The result has limits that the experiments and the theory give numbers for. The guarantee on the long-run frequency of hits is vacuous at the length of a stress window (the bound averages \thAMeanSimple{} percentage points), so the repair is an empirical result; the repair is bounded by the band and by the scale of the model's error (a ratio of \thCCmaxOne{} at the 1\% level for a Student-$t$ model with five degrees of freedom), and the band binds the level well before the fixed point leaves it; a coverage guarantee carries no independence (\cref{prop:independence}); the expected shortfall remains understated (mean $Z_2$ \tenMeanZtwoAci{} wrapped against \tenMeanZtwoUnw{}); the delta-normal model gains no window at the 1\% level; and for volatility-filtered models the wrapper lowers the accuracy of the VaR against the truth in most generator cells. Two other recalibrations that act on the finished VaR pass as many cells (a rolling split-conformal scale, at a cost in the score) or more (a quantile tracker with a larger step, also at a cost), so the wrapper is one of several methods that move coverage, distinguished by keeping the model's own ES at the working level and by a step that has a probabilistic meaning.

The next steps are the tests the claim has not had: portfolios with options and spreads, the models fitted for one level through a refit at the working level, the ES back-tests of the recent literature on the wrapped pair, windows of other markets, and a selection of the step by a rule that does not use the windows on which the wrapper is judged.
